%% GENERATED by figures/build_figures.py from
%%   figures/source/Reviewed_articles_only.xlsx
%% Do not edit by hand: edit the spreadsheet and re-run the script.

\begin{figure*}[pos=!t]
\centering
\begin{tikzpicture}
\begin{axis}[
  jsxbar, bar width=6pt, enlarge y limits=0.05,
  nodes near coords={\pgfmathprintnumber[fixed,precision=0]{\pgfplotspointmeta}},
  title style={align=center, at={(0.5,1)}, anchor=south, yshift=3pt},
  width=0.60\textwidth, height=7.7cm, title={(a) Computing platform named},
  xlabel={Studies}, xmin=0, xmax=71, xtick={0,20,40,60},
  symbolic y coords={NASA Earth Exchange,Microsoft Planetary Computer,Digital Earth Africa,SEPAL FAO,Microsoft Azure,Planet Explorer Portal,Sentinel Hub,ESA CloudToolbox,Hugging Face,Google Cloud Platform,AWS,Google Colab,Local national HPC,Google Earth Engine}, ytick=data, yticklabels={{NASA Earth Exchange},{Microsoft Planetary Computer},{Digital Earth Africa},{SEPAL (FAO)},{Microsoft Azure},{Planet Explorer Portal},{Sentinel Hub},{ESA CloudToolbox},{Hugging Face},{Google Cloud Platform},{AWS},{Google Colab},{Local / national HPC},{Google Earth Engine}},
  yticklabel style={font=\scriptsize},
]
\addplot[draw=none, fill=jsBlue] coordinates {(1,NASA Earth Exchange) (1,Microsoft Planetary Computer) (1,Digital Earth Africa) (1,SEPAL FAO) (1,Microsoft Azure) (1,Planet Explorer Portal) (1,Sentinel Hub) (1,ESA CloudToolbox) (1,Hugging Face) (2,Google Cloud Platform) (3,AWS) (4,Google Colab) (7,Local national HPC) (61,Google Earth Engine)};
\end{axis}
\end{tikzpicture}

\vspace{6pt}
\begin{tikzpicture}
\begin{groupplot}[group style={group size=2 by 1, horizontal sep=32mm}, height=5.2cm]
\nextgroupplot[
  jsxbar, bar width=10pt, enlarge y limits=0.35,
  nodes near coords={\pgfmathprintnumber[fixed,precision=0]{\pgfplotspointmeta}},
  title style={align=center, at={(0.5,1)}, anchor=south, yshift=3pt},
  width=0.38\textwidth, title={(b) Open-science practice},
  xlabel={Studies (of 138)}, xmin=0, xmax=41, xtick={0,10,20,30,40},
  symbolic y coords={Interactive app,Data or code repository,Data link given}, ytick=data, yticklabels={{Interactive app},{Data or code repository},{Data link given}},
]
\addplot[draw=none, fill=jsTeal] coordinates {(5,Interactive app) (25,Data or code repository) (33,Data link given)};
\nextgroupplot[
  jsxbarstack, bar width=6pt, enlarge y limits=0.06,
  title style={align=center, at={(0.5,1)}, anchor=south, yshift=3pt},
  width=0.38\textwidth, title={(c) Repository named, by host},
  xlabel={Studies}, xmin=0, xmax=12, xtick={0,3,6,9},
  symbolic y coords={Source Cooperative,Kaggle,Nakala,GEE,Fordatis record,Mendeley,Figshare,Dataverse,Github,Zenodo}, ytick=data, yticklabels={{Source Cooperative},{Kaggle},{Nakala},{GEE},{Fordatis record},{Mendeley},{Figshare},{Dataverse},{Github},{Zenodo}},
  yticklabel style={font=\scriptsize},
  legend style={at={(0.98,0.04)}, anchor=south east, font=\scriptsize},
]
\addplot[draw=none, fill=jsGreen] coordinates {(1,Source Cooperative) (0,Kaggle) (1,Nakala) (0,GEE) (1,Fordatis record) (1,Mendeley) (2,Figshare) (3,Dataverse) (0,Github) (9,Zenodo)};
\addlegendentry{archival}
\addplot[draw=none, fill=jsRust] coordinates {(0,Source Cooperative) (1,Kaggle) (0,Nakala) (1,GEE) (0,Fordatis record) (0,Mendeley) (0,Figshare) (0,Dataverse) (5,Github) (0,Zenodo)};
\addlegendentry{mutable host}
\end{groupplot}
\end{tikzpicture}
\caption{Computational infrastructure and openness across the 138 reviewed studies.
(a) Computing platforms named by the 69 studies that report one; a study naming
several is counted once per platform. (b) Open-science practice as a count of all
138 studies. (c) The 25 studies naming a data or code repository, by host, with
archival repositories distinguished from mutable code and data hosts.}
\label{fig:infrastructure}
\end{figure*}
